% !TEX root =./ECE444_Practice_main.tex
%
% L8 practice set (Dipole Simulation Lab) -- segmentation arithmetic, sim-vs-theory
% reconciliation, convergence interpretation, the average-gain energy audit, and
% why a method-of-moments code must solve for the current before it can radiate.

\fullwidth{\textbf{Learning Objective 2.2: } I can simulate a dipole antenna using an
EM simulation tool and interpret the results against analytical predictions.
\\[4pt]\normalfont\small Show your work. A \textbf{Documentation} line at the end
is required (write \textbf{None} if you did not collaborate).}

\begin{questions}

%%%%%%%%%%%%%%%%%%%%%%%%%%%%%%%%%%%%%%%%%%%%%%%%%%%%%%%%%%%%%%%%%%%%%%%%%%%%%%%
\question \textbf{Segmentation arithmetic.} You are about to model a straight
center-fed dipole at $f = 2.45 \ghz$ in NEC. The wire is 0.8 mm in diameter,
so the radius is $a = 0.4\ \text{mm}$. Before touching the keyboard, size the
model on paper.

\begin{parts}

	\begin{minipage}[t]{0.58\linewidth}
		\part Find the free-space wavelength and the physical length of a
		half-wave dipole at this frequency.
		\begin{solution}[1.1in]
		\eq{ \lambda &= \frac{c}{f} = \frac{3\ee{8}\ \text{m/s}}{2.45\ee{9}\ \hz} = 0.1224\ \m = 122.4\ \text{mm} \\
		     \frac{\lambda}{2} &= 61.2\ \text{mm} }
		\end{solution}
	\end{minipage}\hfill%
	\begin{minipage}[t]{0.37\linewidth}
		\raggedleft \vspace{12pt}
		$\lambda$ = \ansbox[1.3in]{$122.4\ \text{mm}$} \\[10pt]
		$\lambda/2$ = \ansbox[1.3in]{$61.2\ \text{mm}$}
	\end{minipage}

	\begin{minipage}[t]{0.58\linewidth}
		\part You choose $N = 15$ segments. Find the segment length $\Delta$ and
		show that it satisfies both NEC bounds: $\Delta < \lambda/20$ and
		$\Delta > 8a$.
		\begin{solution}[1.4in]
		\eq{ \Delta &= \frac{61.2\ \text{mm}}{15} = 4.08\ \text{mm} = 0.033\ \lambda }
		Upper bound: $\lambda/20 = 6.12\ \text{mm}$, and $4.08 < 6.12$, so the
		segment is short enough that phase barely changes across it.
		Lower bound: $8a = 8(0.4) = 3.2\ \text{mm}$, and $4.08 > 3.2$, so the
		thin-wire kernel is still valid. Both bounds pass, and 15 segments over a
		half wavelength sits inside the 10--20 rule of thumb.
		\end{solution}
	\end{minipage}\hfill%
	\begin{minipage}[t]{0.37\linewidth}
		\raggedleft \vspace{12pt}
		$\Delta$ = \ansbox[1.3in]{$4.08\ \text{mm}$}
	\end{minipage}

	\begin{minipage}[t]{0.58\linewidth}
		\part What is the largest \emph{odd} segment count this wire will tolerate
		before it violates the $\Delta > 8a$ rule?
		\begin{solution}[1.2in]
		\eq{ \Delta_{min} &= 8a = 3.2\ \text{mm} \\
		     N_{max} &= \frac{61.2\ \text{mm}}{3.2\ \text{mm}} = 19.1 \rightarrow N = 19 }
		This is a fat wire: the refinement ceiling sits barely above the
		segmentation you already chose. If a convergence study demanded 41
		segments, this model could not supply them --- the answer would be to use
		the extended thin-wire kernel or a thinner wire, not more segments.
		\end{solution}
	\end{minipage}\hfill%
	\begin{minipage}[t]{0.37\linewidth}
		\raggedleft \vspace{12pt}
		$N_{max}$ = \ansbox[1.3in]{$19$}
	\end{minipage}

	\part In one or two sentences, explain why the segment count must be odd.
		\begin{solution}[0.9in]
		The voltage source is applied to a single segment, and that segment is the
		antenna's feed point. An odd count puts one segment exactly at the center
		of the wire, so the feed is symmetric. With an even count the center falls
		on a junction between two segments and the source must be placed on one
		side of center, which breaks the symmetry of the model and shifts the
		reported input impedance.
		\end{solution}

\end{parts}

\newpage
%%%%%%%%%%%%%%%%%%%%%%%%%%%%%%%%%%%%%%%%%%%%%%%%%%%%%%%%%%%%%%%%%%%%%%%%%%%%%%%
\question \textbf{Predict, then account for the difference.} You model the
2.45 GHz dipole of Question 1 at \emph{exactly} $\lambda/2 = 61.2\ \text{mm}$,
with 15 segments and a 1 V source on the center segment. NEC returns
$Z_{in} = 86.9 + j46.2 \ohms$ and a peak gain of 2.17 dBi.

\begin{parts}

	\part State the four predictions Lesson 7 makes for a half-wave dipole:
	input impedance at exactly $\lambda/2$, resonant length, input resistance at
	resonance, and directivity in dBi.
		\begin{solution}[1.0in]
		From the sinusoidal-current analysis: $Z_{in} \approx 73 + j42.5 \ohms$ at
		exactly $\lambda/2$; the antenna is resonant slightly short, at about
		$0.47$ to $0.48\ \lambda$; the input resistance there is about $70 \ohms$;
		and the directivity is $D = 1.64$, or $2.15\ \text{dBi}$, with a half-power
		beamwidth of $78\mydeg$.
		\end{solution}

	\begin{minipage}[t]{0.58\linewidth}
		\part Compute the percent difference between the simulated and predicted
		values of $R_{in}$ and of $X_{in}$.
		\begin{solution}[1.3in]
		\eq{ \frac{86.9 - 73}{73} &= 0.190 \rightarrow 19.0\% \\
		     \frac{46.2 - 42.5}{42.5} &= 0.087 \rightarrow 8.7\% }
		The reactance is close; the resistance is not. Note also that the
		simulated gain, 2.17 dBi against a predicted 2.15 dBi, is off by
		1\% --- the pattern quantities agree even though the impedance does not.
		\end{solution}
	\end{minipage}\hfill%
	\begin{minipage}[t]{0.37\linewidth}
		\raggedleft \vspace{12pt}
		$R_{in}$ error = \ansbox[1.2in]{$19.0\%$} \\[10pt]
		$X_{in}$ error = \ansbox[1.2in]{$8.7\%$}
	\end{minipage}

	\part In two or three sentences, account for the resistance difference. Be
	specific about the mechanism --- ``simulation error'' earns nothing.
		\begin{solution}[1.3in]
		The value $73 + j42.5 \ohms$ is the impedance of an assumed sinusoidal
		current, and a sinusoid resonates when the wire is about $0.486\ \lambda$
		long. A real wire has finite radius, stores energy in the near field
		around that radius, and resonates shorter --- near $0.47\ \lambda$. A wire
		cut to exactly $\lambda/2$ is therefore already about 5\% long: it is
		inductive, and its resistance has already climbed well up the steep part
		of the $R_{in}$ versus length curve. The two numbers describe two
		different antennas, not one antenna and one error.
		\end{solution}

	\begin{minipage}[t]{0.58\linewidth}
		\part Trimming brings this wire to resonance at $0.473\ \lambda$. Find the
		resonant length in millimetres and the percent shortening from
		$\lambda/2$.
		\begin{solution}[1.2in]
		\eq{ L_{res} &= 0.473(122.4\ \text{mm}) = 57.9\ \text{mm} \\
		     \text{shortening} &= \frac{61.2 - 57.9}{61.2} = 0.054 \rightarrow 5.4\% }
		Roughly a 5\% trim, which is the standard rule of thumb for bringing a
		dipole to resonance.
		\end{solution}
	\end{minipage}\hfill%
	\begin{minipage}[t]{0.37\linewidth}
		\raggedleft \vspace{12pt}
		$L_{res}$ = \ansbox[1.2in]{$57.9\ \text{mm}$} \\[10pt]
		trim = \ansbox[1.2in]{$5.4\%$}
	\end{minipage}

\end{parts}

\newpage
%%%%%%%%%%%%%%%%%%%%%%%%%%%%%%%%%%%%%%%%%%%%%%%%%%%%%%%%%%%%%%%%%%%%%%%%%%%%%%%
\question \textbf{Reading a convergence study.} A classmate models a dipole of
length exactly $0.5\ \lambda$ with wire radius $a = 0.001\ \lambda$ and re-runs it
at five segment counts:

\begin{center}
\begin{tabular}{c c c}
\hline
\textbf{Segments $N$} & \textbf{$R_{in}$ ($\Omega$)} & \textbf{$X_{in}$ ($\Omega$)} \\
\hline
5  & 79.9 & $+35.9$ \\
11 & 81.9 & $+43.9$ \\
21 & 83.4 & $+45.6$ \\
41 & 84.5 & $+46.5$ \\
81 & 85.4 & $+47.2$ \\
\hline
\end{tabular}
\end{center}

\begin{parts}

	\begin{minipage}[t]{0.58\linewidth}
		\part Compute the percent change in $R_{in}$ from each row to the next.
		\begin{solution}[1.2in]
		\eq{ 5 \rightarrow 11: \quad \frac{81.9-79.9}{79.9} &= 2.5\% \\
		     11 \rightarrow 21: \quad \frac{83.4-81.9}{81.9} &= 1.8\% \\
		     21 \rightarrow 41: \quad \frac{84.5-83.4}{83.4} &= 1.3\% \\
		     41 \rightarrow 81: \quad \frac{85.4-84.5}{84.5} &= 1.1\% }
		The steps shrink, but slowly: $R_{in}$ is still creeping up by about a
		percent per doubling.
		\end{solution}
	\end{minipage}\hfill%
	\begin{minipage}[t]{0.37\linewidth}
		\raggedleft \vspace{12pt}
		last step = \ansbox[1.2in]{$1.1\%$}
	\end{minipage}

	\begin{minipage}[t]{0.58\linewidth}
		\part Find the largest segment count the $\Delta > 8a$ rule permits for
		this wire. Which rows of the table violate it?
		\begin{solution}[1.1in]
		\eq{ \Delta &= \frac{0.5\ \lambda}{N} > 8a = 0.008\ \lambda \\
		     N &< \frac{0.5}{0.008} = 62.5 \rightarrow N_{max} = 61 }
		The $N = 81$ row breaks the rule --- its segments are only about five
		radii long --- so it sits outside the range where the standard thin-wire
		kernel is trustworthy and should not be used as evidence of anything.
		\end{solution}
	\end{minipage}\hfill%
	\begin{minipage}[t]{0.37\linewidth}
		\raggedleft \vspace{12pt}
		$N_{max}$ = \ansbox[1.2in]{$61$}
	\end{minipage}

	\part Which segment count would you defend for this model, and why? Your
	answer must reference both the convergence trend and the kernel limit.
		\begin{solution}[1.2in]
		41 segments. Each further refinement moves $R_{in}$ by only about one
		percent, well below what you could measure on the bench, and
		$\Delta = 0.0122\ \lambda$ clears the $0.008\ \lambda$ floor with room to
		spare for a shorter trimmed wire. Going to 81 buys a percent of drift at
		the cost of leaving the kernel's valid range.
		\end{solution}

	\part Explain in one or two sentences why gain converges at a lower segment
	count than input impedance does.
		\begin{solution}[0.9in]
		Gain is an integral over the entire current, so errors in individual
		segments partly cancel and the total settles quickly. Input impedance is
		read from one segment --- the one carrying the source --- so it inherits
		the full local error of the feed model, which itself changes every time
		the segment length changes.
		\end{solution}

\end{parts}

\newpage
%%%%%%%%%%%%%%%%%%%%%%%%%%%%%%%%%%%%%%%%%%%%%%%%%%%%%%%%%%%%%%%%%%%%%%%%%%%%%%%
\question \textbf{The average-gain audit.} After a full-sphere pattern run of a
lossless dipole in free space, NEC reports an average power gain of 1.15.

\begin{parts}

	\part What does an average power gain of exactly 1.000 tell you about a
	lossless free-space model, and why is that a test rather than a result?
		\begin{solution}[1.1in]
		Average power gain is the gain averaged over the whole sphere, which is
		the total radiated power divided by the total input power. For a lossless
		antenna in free space, every watt delivered to the terminals must leave as
		radiation, so the average must be 1.000. It is a test because you already
		know the answer: any deviation is not physics, it is a defect in the
		model, and it is the cheapest error check available.
		\end{solution}

	\begin{minipage}[t]{0.58\linewidth}
		\part Express the reported 1.15 in decibels, and state whether you would
		accept the run.
		\begin{solution}[1.2in]
		\eq{ 10\ \mylog{1.15} &= 0.61 \db }
		Reject it. The model claims to radiate 15\% more power than it was fed,
		which is impossible; a healthy run lands inside roughly $\pm 5\%$, or
		about $\pm 0.2\ \text{dB}$. Nothing else in this output file --- impedance,
		gain, pattern --- can be trusted until the cause is found.
		\end{solution}
	\end{minipage}\hfill%
	\begin{minipage}[t]{0.37\linewidth}
		\raggedleft \vspace{12pt}
		avg gain = \ansbox[1.1in]{$0.61 \db$}
	\end{minipage}

	\part Give two plausible causes of an average gain this far from 1.000, and
	the specific check that would confirm each.
		\begin{solution}[1.3in]
		First, a segmentation violation: segments shorter than about eight wire
		radii break the thin-wire kernel and corrupt the currents. Check it by
		computing $\Delta$ and $8a$ by hand, and re-run with fewer segments or a
		thinner wire. Second, a geometry or excitation error --- wires that do not
		actually connect, a duplicated wire, or a source placed on a segment that
		is not at the feed. Check it by plotting the geometry and the segment
		numbering, and confirming the source tag and segment index against the
		wire card. A third common cause is requesting the average over less than a
		full sphere, which makes the number meaningless rather than wrong.
		\end{solution}

\end{parts}

\newpage
%%%%%%%%%%%%%%%%%%%%%%%%%%%%%%%%%%%%%%%%%%%%%%%%%%%%%%%%%%%%%%%%%%%%%%%%%%%%%%%
\question \textbf{Why the current comes first.} A classmate claims that NEC
``computes the radiation pattern directly from the geometry.''

\begin{parts}

	\part In two or three sentences, explain why a method-of-moments code must
	solve for the current distribution before it can produce a pattern. Refer to
	the radiation integral.
		\begin{solution}[1.3in]
		The far field is the radiation integral of the current over the antenna:
		the geometry alone tells you where the current can flow, not how much
		flows where or with what phase. The method of moments exists precisely to
		supply that missing current --- it discretizes the wire, enforces
		$E_{tan} = 0$ on the conductor, and solves a linear system for the segment
		currents. Only then does it evaluate the same radiation integral you would
		evaluate by hand, so the geometry sets up the problem and the current
		answers it.
		\end{solution}

	\begin{minipage}[t]{0.58\linewidth}
		\part The solver drives the center segment with $1\ \text{V}$ and reports a
		feed current of $8.90 - j4.80\ \text{mA}$. Find $Z_{in}$, then the magnitude
		of the reflection coefficient and the VSWR referred to $50 \ohms$.
		\begin{solution}[1.6in]
		\eq{ Z_{in} &= \frac{V}{I} = \frac{1}{\lp 8.90 - j4.80 \rp \ee{-3}} \\
		     &= \frac{\lp 8.90 + j4.80 \rp \ee{-3}}{1.0225 \ee{-4}} = 87.0 + j46.9 \ohms \\
		     \Gamma &= \frac{Z_{in} - 50}{Z_{in} + 50} = \frac{37.0 + j46.9}{137.0 + j46.9} \\
		     \vert \Gamma \vert &= \frac{59.7}{144.8} = 0.412 \\
		     \text{VSWR} &= \frac{1 + 0.412}{1 - 0.412} = 2.40 }
		\end{solution}
	\end{minipage}\hfill%
	\begin{minipage}[t]{0.37\linewidth}
		\raggedleft \vspace{12pt}
		$Z_{in}$ = \ansbox[1.4in]{$87.0 + j46.9 \ohms$} \\[10pt]
		$\vert \Gamma \vert$ = \ansbox[1.2in]{$0.412$} \\[10pt]
		VSWR = \ansbox[1.2in]{$2.40$}
	\end{minipage}

	\part You need this antenna matched to $50 \ohms$. Would you trim the wire to
	resonance or add a matching network? Give the trade-off in one or two
	sentences.
		\begin{solution}[1.1in]
		Trim first. Trimming costs nothing, removes the $+j46.9 \ohms$ entirely,
		and leaves about $72 \ohms$ real, which is already a VSWR near 1.4 on a
		$50 \ohms$ line --- acceptable for most links. A matching network would
		bring the residual mismatch closer to unity but adds parts, loss, and a
		bandwidth restriction, so it is worth doing only when the remaining
		mismatch actually costs you something in the link budget.
		\end{solution}

\end{parts}

\vspace{1.5em}
\noindent\textbf{Documentation:} \rule[-2pt]{0.6\linewidth}{0.4pt}

\end{questions}
